\section{Related Works}

\textbf{Optimizers for Deep Learning.} A long line of work has studied optimization
for deep learning, incorporating insights from classical optimization literature~\cite{robbins1951stochastic,nesterov1983method,duchi2011adaptive} and domain knowledge about deep neural networks~\cite{pmlr-v28-sutskever13}. (i) Early insights motivated the rise of optimizers that use adaptive learning rates based on second-order momentum~\cite{tieleman2012lecture,zeiler2012adadeltaadaptivelearningrate,kingma2017adammethodstochasticoptimization}. Adam~\cite{kingma2017adammethodstochasticoptimization} later became the default baseline for optimizers with adaptive learning rates. Since then, improvements over Adam and SGD have been proposed, with notable examples including Nesterov Adam~\cite{dozat2016incorporating} and AdamW with decoupled weight decay~\cite{loshchilov2019decoupled}. Other improvements include addressing the convergence of Adam on convex loss~\cite{j.2018on,zaheer2018adaptive,taniguchi2024adoptmodifiedadamconverge}, considering interpolation between Adam and SGD to improve generalization~\cite{luo2019adaptivegradientmethodsdynamic,xie2022adaptiveinertiadisentanglingeffects}, performing further variance reductions on optimizer updates~\cite{liu2021varianceadaptivelearningrate,zhang2019lookaheadoptimizerksteps,yuan2025mars,xie2024adanadaptivenesterovmomentum,pagliardini2024ademamixoptimizerbetterfaster,zhuang2020adabeliefoptimizeradaptingstepsizes}, incorporating momentum on weights~\cite{ivgi2023dogsgdsbestfriend,defazio2024roadscheduled}, allowing easier hyperparameter tuning~\cite{pmlr-v202-defazio23a,mishchenko2024prodigyexpeditiouslyadaptiveparameterfree,defazio2024optimallineardecaylearning}, reducing memory usage by incorporating the structure of 
neural networks~\cite{shazeer2018adafactoradaptivelearningrates,zhang2025adamminiusefewerlearning,zhu2025apollosgdlikememoryadamwlevel,luo2023cameconfidenceguidedadaptivememory,modoranu2024microadam,zhao2024galorememoryefficientllmtraining}, and modifying the algorithm to allow larger batch sizes~\cite{you2017largebatchtrainingconvolutional,you2020largebatchoptimizationdeep}.
(ii) Starting from Preconditioned SGD~\cite{Li_2018} and Shampoo~\cite{gupta2018shampoopreconditionedstochastictensor}, another line of optimizer design~\cite{morwani2024newperspectiveshampoospreconditioner,eschenhagen2023kronecker,martens2020optimizingneuralnetworkskroneckerfactored,li2018preconditionermatrixliegroup,eschenhagen2025purifyingshampooinvestigatingshampoos} began to incorporate matrix preconditioners instead of simple scalar preconditioners. Techniques including learning rate grafting~\cite{agarwal2020disentanglingadaptivegradientmethods}, blocking, and distributed methodology~\cite{anil2021scalablesecondorderoptimization} have since been proposed. These matrix-based approaches later led to the theory of modular duality in deep learning optimization~\cite{bernstein2024modulardualitydeeplearning,bernstein2024oldoptimizernewnorm,large2024scalableoptimizationmodularnorm} and new optimizers including Muon~\cite{jordan2024muon} and Scion~\cite{pethick2025trainingdeeplearningmodels}.
(iii) Motivated by Newton's algorithm, there has been a line of work that tries to incorporate Hessian information~\cite{becker:improving,Yao_Gholami_Shen_Mustafa_Keutzer_Mahoney_2021,schaul2013peskylearningrates,Yao_Gholami_Shen_Mustafa_Keutzer_Mahoney_2021}. (iv) Symbolic discovery of optimizers~\cite{chen2023symbolic} has discovered a memory-efficient SignGD~\cite{pmlr-v80-bernstein18a} variant called Lion, which claims to outperform Adam on a wide range of tasks.

\textbf{Optimization for Pretraining.} Since~\cite{brown2020languagemodelsfewshotlearners}, the cost of pretraining has increased dramatically. One of the challenges is how to choose hyperparameters with minimal cost.  
A pivotal line of work in this direction is the tensor program series that allows for extrapolating some hyperparameters across scales~\cite{yang2020tensor,yang2020tensor2,yang2020scalinglimitswideneural,yang2021tensorprogramsiwide,yang2023tensorprogramsivbadaptive,yang2022tensorprogramsvtuning,yang2024spectralconditionfeaturelearning}. Empirical results suggest that fitting a power law to scale hyperparameters is now common practice for large models~\cite{li2025predictablescalei,everett2024scalingexponentsparameterizationsoptimizers,liu2025muonscalablellmtraining,deepseekai2024deepseekllmscalingopensource,zhang2025doescriticalbatchsize}. We have also incorporated this approach in our paper. 
Another popular line of research is designing better optimizers specifically for pretraining. These optimizers are our main objects of study. An (incomplete) list of optimizers and their claimed speedups over AdamW includes Sophia~\cite{liu2024sophiascalablestochasticsecondorder} (2$\times$), Soap~\cite{vyas2025soapimprovingstabilizingshampoo} (1.4$\times$), Muon~\cite{jordan2024muon,liu2025muonscalablellmtraining} (2$\times$), MARS~\cite{yuan2025mars} (2$\times$), Cautious AdamW~\cite{liang2025cautious} (2$\times$), Block-wise Learning Rate Adam~\cite{wang2025sharpnessdisparityprincipletransformers} (2$\times$), FOCUS~\cite{liu2025focusorderconcentratedupdating} (2$\times$), SWAN~\cite{ma2025swansgdnormalizationwhitening} (2$\times$), DION~\cite{ahn2025diondistributedorthonormalizedupdates} (3$\times$), and SPlus~\cite{frans2025stablewhiteningoptimizerefficient} (2$\times$). We present a comparison of setups with these prior works in ~\Cref{sec:comparison-with-prior}.

\textbf{Re-evaluation Methodology.} Our work is a rigorous evaluation of optimizers for pretraining. Rigorous evaluation has been an important part of deep learning research to clarify the current status of research and move the community forward. \cite{jiang2019fantasticgeneralizationmeasures} critically examines metrics for predicting LLMs' generalization capability and has facilitated research on understanding loss landscape sharpness and generalization and new optimizers such as SAM~\cite{foret2021sharpnessawareminimizationefficientlyimproving}. \cite{schmidt2021descendingcrowdedvalley} re-evaluated optimizers at that time and showed that (i) which optimizer is optimal is problem-specific, and (ii) rigorous hyperparameter tuning is required and unequal tuning can account for most of the claimed speedup. Unlike~\cite{schmidt2021descendingcrowdedvalley}, we evaluate the LLM pretraining task and include modern optimizers they did not test. However, we arrive at a similar conclusion: rigorous and fair hyperparameter tuning is still not the norm, but rather the exception in optimizer design research.~\cite{zhao2025deconstructingmakesgoodoptimizer} also examines how different optimizers perform on the pretraining tasks. However, the focus of~\cite{zhao2025deconstructingmakesgoodoptimizer} is on understanding loss structure, and the optimizers tested in the paper are shown to have slower convergence compared to AdamW, whereas the optimizers tested in this paper show small but significant improvements in convergence speed. The Algoperf competition~\cite{kasimbeg2025acceleratingneuralnetworktraining} evaluates different optimizers across different settings and arrives at a similar conclusion that (i) matrix-based optimizers and (ii) denoising methods such as Nesterov momentum lead to speedup over AdamW. Our works focus on the pretraining setting and investigates the effect of scaling dataset and model sizes.~\citet{kaddour2023traingainrevisitingefficient} compares optimizers including Sophia, as well as some other algorithmic improvements, in the setting of training Bert and T5-like models. They also observe the ordering of optimizers can flip during learning rate decay, and AdamW remains competitive when properly tuned.

\paragraph{Comparison with concurrent work~\citet{semenov2025benchmarkingoptimizerslargelanguage}} \citet{semenov2025benchmarkingoptimizerslargelanguage} investigates the same topic of benchmarking optimizers for pre-training. These two works agree on many high-level points, e.g., (i) non-zero weight decay and decaying to a small learning rate are essential for pretraining; (ii) variance-reduced AdamW variants such as Mars show non-trivial speedups over vanilla AdamW. 

However, our results differ in the relative performance of matrix-level optimizers. Our paper shows that matrix-level optimizers such as Muon can achieve significant speedups over variance-reduced versions of AdamW, such as Mars, their study finds that AdEMaMix and Mars outperform Muon.  Our initial investigation suggests this discrepancy largely stems from the differences in the batch sizes used in experiments. Their most extensively tuned experiments on 130M models use a batch size of only 0.1M and 0.02M tokens, whereas our experiments operate with tuned batch sizes that are not smaller than 0.4M tokens. This difference is likely a result of different hardware regimes. We leverage 128 TPU-v5lite chips (which are approximately equivalent to 12 H100 GPUs), where only larger batches (larger than 0.4M) can fully utilize the parallelism in the compute. On the other side, their experiments appear to be primarily conducted on 1-8 H100 GPUs, where smaller batches might be preferable. Since Mars and AdEMAMix both perform gradient averaging and variance reduction, these methods are advantangeous in their noise-dominated small-batch regime, whereas in our larger-batch setting these benefits diminish and matrix-level optimizers become more competitive. 

In their larger scale experiments,~\citet{semenov2025benchmarkingoptimizerslargelanguage} increase both model size to 720M parameters and batch size to 1M tokens, respectively, which is closer to our setting. However, we differ in hyperparameter tuning methodology: (i) we conduct more extensive sweeps over learning rates on our 520M scales experiments and generally find higher values (4e-3 to 8e-3) compared to their choices (1e-3 to 2e-3), which are transferred from smaller scale and for some optimizers are not tuned; and (ii) for Muon, we separately tune the learning rates for the embedding layers and the matrix weights, which also improves its performance. These differences highlight the sensitivity of optimizer benchmarking to hardware and tuning strategies, underscoring the importance of carefully controlled experimental design when comparing optimizer performance.
